\chapter{Program Structure}
\label{ch:structure}

This chapter describes the robot's hardware, how the code is split into files,
and the build step that combines those files before they are sent to the robot.
Understanding this structure first makes the later chapters easier to follow,
because each later chapter focuses on one of these files.

\section{Hardware}
\label{sec:hardware}

The V5 Brain has 21 numbered \emph{smart ports}. Each motor or sensor plugs into
one port, and the program must be told which port each device is on.
\tref{tab:ports} lists the devices.

\begin{table}[h]
  \centering
  \begin{tabular}{@{}cll@{}}
    \toprule
    Port & Device & Cartridge \\
    \midrule
    1  & Right drive motor & green, 18:1 \\
    2  & Left drive motor  & green, 18:1 \\
    4  & AI Vision sensor  & -- \\
    10 & Arm motor         & red, 36:1 \\
    15 & Claw motor        & red, 36:1 \\
    9  & VEXnet radio      & -- \\
    \bottomrule
  \end{tabular}
  \caption{Devices and the ports they are plugged into.}
  \label{tab:ports}
\end{table}

A \emph{cartridge} is the gearbox inside a V5 motor. A green 18:1 cartridge
gives a top speed of 200 revolutions per minute (rpm), which suits driving. A
red 36:1 cartridge gives 100 rpm but twice the torque (turning force), which
suits an arm that has to lift weight. The program must declare the correct
cartridge for each motor, because the motor uses it to convert its internal
rotation into output degrees. We return to what happens when this is wrong in
\sref{sec:gotchas}.

The radio on port 9 is managed by the Brain's own firmware, so the program
never declares it.

Note: in the Python SDK, ports are numbered from zero, so \code{Ports.PORT1} has
the value 0. This is why the port scanner in \sref{sec:portscan} adds 1 before
displaying a port number.

\section{File Layout}
\label{sec:layout}

The program is split into five Python files inside \code{src/}, plus one build
script. \tref{tab:files} summarises them.

\begin{table}[h]
  \centering
  \begin{tabularx}{\textwidth}{@{}lX@{}}
    \toprule
    File & Contents \\
    \midrule
    \code{config.py} & Every device and every tunable number. \\
    \code{vision.py} & Cone recognition and AI Vision diagnostics. \\
    \code{autonomousroutine.py} & Movement maths, calibration tools, and \code{autonomous()}. \\
    \code{drivercontrol.py} & Stick processing, the arm and claw, cone assist, and \code{user_control()}. \\
    \code{main.py} & Chooses which mode runs. \\
    \code{build.py} & Combines the five files into \code{src/download.py}. \\
    \bottomrule
  \end{tabularx}
  \caption{The files in the project and what each one holds.}
  \label{tab:files}
\end{table}

We keep every tunable number in \code{config.py}. This is useful because tuning
the robot usually means changing a speed, a distance, or a tolerance, and
having them in one place means we never have to search through the logic to
find them. Each constant has a short comment giving its unit or meaning, for
example:

\lstinputlisting[firstline=35,firstnumber=35,lastline=45,caption={Autonomous constants in \code{src/config.py}, lines 35--45.},label=lst:autoconst]{\src{config.py}}

\section{The Single-File Constraint}
\label{sec:singlefile}

The V5 Brain can only run \textbf{one} Python file. The project's settings are
stored in \code{vex_project_settings.json}, inside the \code{.vscode} folder.
This file has a single \code{main} entry that names the file to download, and
there is no way to list several files. This
creates a problem: a line such as \code{from config import *} works in the
editor, because the editor can see \code{config.py}, but fails on the robot with
an \code{ImportError}, because \code{config.py} was never sent.

Before 22 September, the whole program lived in one 425-line \code{main.py}.
This worked, but it was difficult to read and edit because driving, autonomous
and settings were all mixed together. Splitting it into files (commit
\commit{0607385}) made each part easier to find, but only because we also wrote
a build step to undo the split before download.

\section{The Build Step}
\label{sec:build}

\code{build.py} reads the five files in a fixed order and joins them into one
file, \code{src/download.py}. The project is configured to download that file
instead of \code{main.py}. \fref{fig:build} shows the process.

\begin{figure}[h]
  \centering
  \begin{tikzpicture}[
      file/.style={draw, rounded corners, minimum width=3.4cm, minimum height=0.75cm, font=\ttfamily\small},
      outfile/.style={draw, thick, rounded corners, minimum width=3.4cm, minimum height=0.9cm, font=\ttfamily\small, fill=codebg},
      >={Stealth}]
    \node[file] (c) at (0,  2.4) {config.py};
    \node[file] (v) at (0,  1.2) {vision.py};
    \node[file] (a) at (0,  0.0) {autonomousroutine.py};
    \node[file] (d) at (0, -1.2) {drivercontrol.py};
    \node[file] (m) at (0, -2.4) {main.py};
    \node[draw, align=center, minimum width=3.2cm, minimum height=1.6cm, font=\small] (b) at (5,0)
      {\texttt{build.py}\\[2pt] remove local imports\\ check it compiles};
    \node[outfile] (o) at (10,0) {download.py};
    \node[font=\small, below=0.2cm of o] {sent to the Brain};
    \foreach \n in {c,v,a,d,m} \draw[->] (\n.east) -- (b.west);
    \draw[->, thick] (b) -- (o);
  \end{tikzpicture}
  \caption{\code{build.py} joins the five source files, in dependency order, into the single file the Brain runs (self-made).}
  \label{fig:build}
\end{figure}

The order matters. Python reads a file from top to bottom, so a function must be
defined before anything that uses it runs. Since every other file reads values
from \code{config.py}, it must come first. The order is stored in a list:

\lstinputlisting[firstline=12,firstnumber=12,lastline=19,caption={Module order and import patterns in \code{build.py}, lines 12--19.},label=lst:buildorder]{../build.py}

The two regular expressions (text patterns) match the import lines that would
break on the robot. \code{LOCAL_IMPORT} matches imports of our own files, which
are no longer needed once everything is in one file. \code{VEX_IMPORT} matches
\code{from vex import *}, which every file repeats; the build keeps just one
copy at the top. The main loop then removes those lines from each file and
joins what remains:

\lstinputlisting[firstline=36,firstnumber=36,lastline=59,caption={Flattening and checking in \code{build.py}, lines 36--59.},label=lst:buildloop]{../build.py}

Line 54 is a safety check. \code{compile()} asks Python to read the combined
file without running it. If the combination has a syntax error, the build stops
with a clear message instead of writing a broken file. Without this check, the
error would only appear on the robot, with a line number pointing into the
nearly 900-line generated file, which is much harder to trace back to the original
source.

Note: \code{download.py} is generated, so it must never be edited by hand. Any
edit would be overwritten the next time \code{build.py} runs. We still commit it
to git, so that a fresh copy of the project can be downloaded to a robot
immediately at a competition, even without Python installed.

\section{Selecting a Mode}
\label{sec:modes}

At a real competition, the field control system tells the robot when to run
autonomous and when to hand over to the driver. During practice there is no
field control, so \code{main()} chooses one mode by hand:

\lstinputlisting[firstline=13,firstnumber=13,lastline=23,caption={Mode selection in \code{src/main.py}, lines 13--23.},label=lst:main]{\src{main.py}}

Exactly one line is left uncommented. The diagnostic modes
(\code{scan_ports}, \code{vision_report}, \code{cone_detector}) and the
calibration mode (\code{turn_tuner}) are kept here so that any of them can be
run on the robot by changing one line, rather than by writing a separate
program.

\section{Summary}

\begin{enumerate}
  \item The robot has two green-cartridge drive motors, two red-cartridge
    manipulator motors and an AI Vision sensor (\sref{sec:hardware}).
  \item The code is split into five files, with every tunable number kept in
    \code{config.py} (\sref{sec:layout}).
  \item The Brain runs only one file, so \code{build.py} joins the five files
    and checks the result compiles before it is downloaded (\sref{sec:build}).
\end{enumerate}

We now explain how the robot moves during the autonomous period.
